\documentclass{article}
\usepackage[utf8]{inputenc}
\usepackage{amsmath,amssymb}
\usepackage[most]{tcolorbox}
\usepackage{geometry}
\geometry{margin=2.5cm}

\newcommand{\step}[1]{\par\noindent\hangindent=3.5em\hangafter=1%
  \makebox[3.5em][l]{\textbf{#1.}}}
\newcommand{\steptag}[1]{\hfill{\small\textsf{[#1]}}}

\newtcolorbox{proofbox}[1]{
  colback=gray!5, colframe=gray!60,
  fonttitle=\bfseries, title={#1},
  breakable, enhanced,
  left=4pt, right=4pt, top=4pt, bottom=4pt,
  before skip=10pt, after skip=10pt
}

\begin{document}

\begin{proofbox}{Fix the def-maincb-witness-ledger datum W supplied by lem...}
\step{1} Fix the def-maincb-witness-ledger datum W supplied by lem-maincb-reset-constant-ledger; every finite-dimensional extended epsilon-C*-algebra A with 0 <= epsilon <= W.epsilon\_MAIN admits a finite-dimensional C*-algebra B=oplus\_C M\_\{|C|\} and an extended W.c0\_cb*W.K\_call*epsilon-isomorphism v:B->A satisfying ||v(I\_B)-I\_A|| <= W.c0\_cb*W.K\_call*epsilon; hence C\_struct=W.c0\_cb*W.K\_call and e\_struct=W.epsilon\_MAIN are finite positive universal witnesses. \steptag{validated} \\[6pt]
\step{1.1} For every finite-dimensional extended epsilon-C*-algebra A with 0 <= epsilon <= W.epsilon\_MAIN, there exist B=oplus\_C M\_\{|C|\} and an extended W.c0\_cb*W.K\_call*epsilon-isomorphism v:B->A with ||v(I\_B)-I\_A|| <= W.c0\_cb*W.K\_call*epsilon. \steptag{validated} \\[6pt]
\step{1.1.1} By lem-maincb-maximal-reset-selection, there exist a positive integer m of maximum source dimension and an extended W.c0\_cb*epsilon-inclusion w:C\textasciicircum{}m->A satisfying ||w(I\_\{C\textasciicircum{}m\})-I\_A|| <= W.c0\_cb*epsilon. \steptag{validated} \\[4pt]
\step{1.1.2} If m>=1 is of maximum source dimension and w:C\textasciicircum{}m->A is an extended W.c0\_cb*epsilon-inclusion with ||w(I\_\{C\textasciicircum{}m\})-I\_A|| <= W.c0\_cb*epsilon, then there exist B=oplus\_C M\_\{|C|\} and an extended W.c0\_cb*W.K\_call*epsilon-isomorphism v:B->A with ||v(I\_B)-I\_A|| <= W.c0\_cb*W.K\_call*epsilon. \steptag{validated} \\[6pt]
\step{1.1.2.1} Put the nonempty set J=\{1,...,m\}, let \{e\_j\} be the projection basis of C\textasciicircum{}m, set P\_j=w(e\_j), R=w(I\_\{C\textasciicircum{}m\}), and t=W.K\_call*epsilon. Then A is an extended t-C*-algebra, every P\_j and R is a t-projection, ||R-I\_A|| <= W.c0\_cb*epsilon <= t, every P\_j is one-dimensional, and j\textasciitilde{}k iff dim S\textasciicircum{}A\_\{P\_j,P\_k\}=1 is an equivalence relation with a finite nonempty class family C. \steptag{validated} \\[6pt]
\step{1.1.2.1.1} By lem-maincb-witness-arithmetic, W.K\_call>=max\{1,W.c0\_cb\}; hence 0<=epsilon<=t and W.c0\_cb*epsilon<=t. The definitions def-extended-epsilon-cstar-algebra and def-epsilon-cstar-algebra show directly, by weakening every epsilon-bound to the larger t-bound at every amplification, that A is extended t-C*. By def-projection-basis, e\_j and I are self-adjoint idempotents of norm one; since w is an extended W.c0\_cb*epsilon-inclusion, def-extended-delta-inclusion and GT-kitaev-def-delta-homomorphism at level one imply P\_j\textasciicircum{}dagger=P\_j, R\textasciicircum{}dagger=R, ||P\_j\textasciicircum{}2-P\_j||<=W.c0\_cb*epsilon<=t, and ||R\textasciicircum{}2-R||<=W.c0\_cb*epsilon<=t. Thus def-delta-projection makes P\_j,R t-projections, and the assumed near-unit bound gives ||R-I\_A||<=W.c0\_cb*epsilon<=t. \steptag{validated} \\[4pt]
\step{1.1.2.1.2} Use the particular e\_sim witness from lem-maincb-corner-equivalence that was also used, together with the particular e\_full witness, to instantiate the single fixed W through lem-maincb-reset-constant-ledger. Under node 1.1.2 and assuming A is extended t-C* and each P\_j is a t-projection, lem-maincb-stage1-maximality gives dim S\textasciicircum{}A\_\{P\_j\}=1 for every j; with def-one-dimensional-delta-projection this makes every P\_j one-dimensional. Lem-maincb-structural-domain-ledger gives t=W.K\_call*epsilon<=e\_sim for that same witness, so lem-maincb-corner-equivalence makes j\textasciitilde{}k iff dim S\textasciicircum{}A\_\{P\_j,P\_k\}=1 an equivalence relation. Since nonempty J is finite, its equivalence classes form a finite nonempty class family C. \steptag{validated} \\[4pt]
\step{1.1.2.2} If nonempty J=\{1,...,m\}, \{e\_j\} is the projection basis of C\textasciicircum{}m, P\_j=w(e\_j), R=w(I\_\{C\textasciicircum{}m\}), and t=W.K\_call*epsilon; if A is extended t-C*, P\_j and R are t-projections, ||R-I\_A|| <= W.c0\_cb*epsilon <= t, all P\_j are one-dimensional, and \textasciitilde{} is an equivalence relation with finite nonempty class family C, then there exist B=oplus\_\{C in C\} M\_\{|C|\} and an extended W.c0\_cb*W.K\_call*epsilon-isomorphism v:B->A satisfying ||v(I\_B)-I\_A|| <= W.c0\_cb*W.K\_call*epsilon. \steptag{validated} \\[6pt]
\step{1.1.2.2.1} Let P\_J=sum\_\{j in J\}P\_j=R and A\_J=S\textasciicircum{}A\_R. Use the particular e\_full witness from lem-maincb-full-corner-identification that was used, together with the particular e\_sim witness, to instantiate the single fixed W through lem-maincb-reset-constant-ledger. Lem-maincb-structural-domain-ledger gives t<=e\_full for that same witness. Since R is a t-projection and ||R-I\_A||<=t in the extended t-C*-algebra A, lem-maincb-full-corner-identification gives Co\_R=Id\_A and S\textasciicircum{}A\_R=A. Hence A\_J=A; and by def-compressed-corner its compressed unit is u\_\{A\_J\}=Co\_R(R)=R. \steptag{validated} \\[4pt]
\step{1.1.2.2.2} Using A\_J=A and u\_\{A\_J\}=R, the class family C supports a MAIN partition state and per-class reset isomorphisms, whose finite recombination gives the asserted B and v with the target map-defect and unit bounds. \steptag{validated} \\[6pt]
\step{1.1.2.2.2.1} Choose the current subset U=J. Since w is an extended W.c0\_cb*epsilon-inclusion and 0<=W.c0\_cb*epsilon<=t, lem-maincb-extended-inclusion-monotone makes w an extended t-inclusion. Using A\_J=A, define a def-maincb-reset-state for J by epsilon\_J=t, B\_J=C\textasciicircum{}m, v\_J=w, its fixed amplification family, d\_J=t, and the extended-inclusion tag. Together with the already fixed finite J, A, w, P\_j, equivalence relation and class family C, this supplies explicitly the def-maincb-partition-state for the same displayed A,w, with current subset J and the stated reset-state reference. \steptag{validated} \\[6pt]
\step{1.1.2.2.2.1.1} From the dependency A\_J=A and u\_\{A\_J\}=R, and from 0<=W.c0\_cb*epsilon<=t, lem-maincb-extended-inclusion-monotone upgrades w to an extended t-inclusion with codomain A\_J. Taking U=J, epsilon\_J=t, B\_J=C\textasciicircum{}m, v\_J=w, the fixed amplifications, d\_J=t, and the extended-inclusion tag gives every field required by def-maincb-reset-state; J is a nonempty union of all classes. Adjoining this reset-state reference to the explicitly fixed finite J,A,w,P\_j, relation and class family gives every field required by def-maincb-partition-state for the same A,w. \steptag{validated} \\[4pt]
\step{1.1.2.2.2.2} For every class C in the finite class family, apply lem-maincb-one-class-extension to this same explicit MAIN partition state and the original w. Its hypotheses are exactly the extended W.c0\_cb*epsilon-inclusion and near-unit bound on w, 0<=epsilon<=W.epsilon\_MAIN, the one-dimensional atomic images, and that C is one equivalence class. It yields a current reset isomorphism v\_C:M\_\{|C$|\}-\rangle$A\_C with A\_C extended epsilon\_C-C*, epsilon\_C<=W.L*epsilon<=W.K\_call*epsilon, d\_C<=W.c0\_cb*epsilon\_C, and ||v\_C(I)-u\_\{A\_C\}||<=W.c0\_cb*epsilon\_C. \steptag{validated} \\[4pt]
\step{1.1.2.2.2.3} Enumerate the finite class family as C\_1,...,C\_q. Applying lem-maincb-stage3-finite-recombination to the explicit partition state and the per-class reset isomorphisms gives B=oplus\_\{a=1\}\textasciicircum{}q M\_\{|C\_a|\} and a current reset isomorphism v:B->A\_J with recorded epsilon\_J\textasciicircum{}out<=W.L*epsilon, map defect d\_J\textasciicircum{}out<=W.c0\_cb*epsilon\_J\textasciicircum{}out, and ||v(I\_B)-u\_\{A\_J\}||<=W.c0\_cb*epsilon\_J\textasciicircum{}out. \steptag{validated} \\[6pt]
\step{1.1.2.2.2.3.1} Because the finite nonempty family C is the family of equivalence classes of J, choose an enumeration C\_1,...,C\_q with q>=1 and union\_\{a=1\}\textasciicircum{}q C\_a=J. Validated node 1.1.2.2.2.1.1 supplies, for the same displayed A,w, the explicit def-maincb-partition-state with this class family, and validated node 1.1.2.2.2.2 supplies for every C\_a an initial current reset isomorphism v\_\{C\_a\}:M\_\{|C\_a$|\}-\rangle$A\_\{C\_a\} whose recorded ambient field epsilon\_\{C\_a\} satisfies epsilon\_\{C\_a\}<=W.L*epsilon, with d\_\{C\_a\}<=W.c0\_cb*epsilon\_\{C\_a\} and ||v\_\{C\_a\}(I)-u\_\{A\_\{C\_a\}\}||<=W.c0\_cb*epsilon\_\{C\_a\}. These are exactly the hypotheses of lem-maincb-stage3-finite-recombination (together with the inherited finite-dimensional extended epsilon-C*-algebra A, 0<=epsilon<=W.epsilon\_MAIN, the original extended W.c0\_cb*epsilon-inclusion and near-unit bounds on w, and the one-dimensional atomic images). Applying that lemma yields a current reset isomorphism v:oplus\_\{a=1\}\textasciicircum{}q M\_\{|C\_a$|\}-\rangle$A\_\{union\_a C\_a\}=A\_J with recorded epsilon\_J\textasciicircum{}out<=W.L*epsilon, d\_J\textasciicircum{}out<=W.c0\_cb*epsilon\_J\textasciicircum{}out, and ||v(I)-u\_\{A\_J\}||<=W.c0\_cb*epsilon\_J\textasciicircum{}out, which is the asserted conclusion. \steptag{validated} \\[4pt]
\step{1.1.2.2.2.4} Put D=W.c0\_cb*W.K\_call*epsilon. The direct sum B=oplus\_\{C in C\}M\_\{|C|\} is a finite-dimensional C*-algebra, and the recombined current reset state tags v:B->A\_J as an extended d\_J\textasciicircum{}out-isomorphism into the compressed algebra A\_J with unit u\_\{A\_J\}=R. The unit-bearing codomain (A,I\_A) is different, so lem-maincb-extended-inclusion-monotone is not applied across these units. Instead use the validated direct check in node 1.1.2.2.2.4.1: node 1.1.2.2.1 gives Co\_R=Id\_A and A\_J=A, hence the multiplication, involution, and amplified matrix norms agree; every non-unit amplified clause weakens from d\_J\textasciicircum{}out to D because 0<=d\_J\textasciicircum{}out<=W.c0\_cb*epsilon\_J\textasciicircum{}out<=W.c0\_cb*W.L*epsilon<=D, while the amplified unit clause follows from ||(I\_n tensor v)(I\_n tensor I\_B)-I\_n tensor I\_A||\_n<=d\_J\textasciicircum{}out+||R-I\_A||<=W.c0\_cb*(W.L+1)*epsilon<=D, using ||R-I\_A||<=W.c0\_cb*epsilon and W.K\_call>=W.L+1. Bijectivity is unchanged under A\_J=A. Thus v:B->(A,I\_A) is an extended D-isomorphism by def-extended-delta-inclusion, and at n=1 the same unit estimate gives ||v(I\_B)-I\_A||<=D. \steptag{validated} \\[6pt]
\step{1.1.2.2.2.4.1} Put D=W.c0\_cb*W.K\_call*epsilon and q=||R-I\_A||. The Stage-3 tag makes v an extended d\_J\textasciicircum{}out-isomorphism into the compressed algebra A\_J, whose unit is R, so it cannot be transferred to the unit-bearing algebra (A,I\_A) by lem-maincb-extended-inclusion-monotone. Instead, Co\_R=Id\_A and A\_J=A imply that the multiplication, involution, and matrix norms of A\_J are those of A; only the distinguished unit changes from R to I\_A. At every amplification the non-unit clauses of the extended d\_J\textasciicircum{}out-isomorphism therefore remain valid into A and weaken to defect D because d\_J\textasciicircum{}out<=W.c0\_cb*epsilon\_J\textasciicircum{}out<=W.c0\_cb*W.L*epsilon<=D. The unit clause changes by at most q: ||(I\_n tensor v)(I\_n tensor I\_B)-I\_n tensor I\_A||\_n <= d\_J\textasciicircum{}out+q <= W.c0\_cb*(W.L+1)*epsilon <= D, using q<=W.c0\_cb*epsilon and W.K\_call>=W.L+1. Thus direct verification of def-extended-delta-inclusion, not monotonicity across different units, makes v:B->(A,I\_A) an extended D-isomorphism and also gives ||v(I\_B)-I\_A||<=D. \steptag{validated} \\[6pt]
\step{1.1.2.2.2.4.1.1} For each n>=1, identify M\_n tensor A\_J with M\_n tensor A as a normed involutive algebra. This is legitimate because node 1.1.2.2.1 gives A\_J=A and Co\_R=Id\_A at every amplification, while def-compressed-corner defines the corner product by compression; hence the corner product becomes the ambient product, and the inherited involution and matrix norms are unchanged. The Stage-3 extended d\_J\textasciicircum{}out-isomorphism tag then supplies linearity, bijectivity, involution preservation, the d\_J\textasciicircum{}out multiplication estimate, and the two-sided (1 plus-or-minus d\_J\textasciicircum{}out) norm estimates for I\_n tensor v with codomain M\_n tensor A; these clauses do not mention the target unit and, since 0<=d\_J\textasciicircum{}out<=D, each implies its D-version. For the sole unit-dependent clause, the corner tag gives ||(I\_n tensor v)(I\_n tensor I\_B)-I\_n tensor R||\_n<=d\_J\textasciicircum{}out. The operator-space direct-sum axiom gives ||I\_n tensor (R-I\_A)||\_n=||R-I\_A||=q, so the triangle inequality gives unit error at most d\_J\textasciicircum{}out+q. From epsilon\_J\textasciicircum{}out<=W.L*epsilon, d\_J\textasciicircum{}out<=W.c0\_cb*epsilon\_J\textasciicircum{}out, q<=W.c0\_cb*epsilon, and W.K\_call>=W.L+1, we get d\_J\textasciicircum{}out+q<=W.c0\_cb*(W.L+1)*epsilon<=D. Thus every amplification is a D-isomorphism into (M\_n tensor A,I\_n tensor I\_A), proving that v is an extended D-isomorphism; n=1 also yields ||v(I\_B)-I\_A||<=D. \steptag{validated} \\[4pt]
\step{1.2} The scalars C\_struct=W.c0\_cb*W.K\_call and e\_struct=W.epsilon\_MAIN are finite, positive, universal, and independent of dimension, amplification, block data, class count, and stage index. \steptag{validated} \\[6pt]
\step{1.2.1} By lem-maincb-reset-constant-ledger and lem-maincb-witness-arithmetic, every field of the fixed ledger W, in particular W.c0\_cb, W.K\_call, and W.epsilon\_MAIN, is positive, finite, universal, and independent of dimension, amplification, block data, class count, and stage index; therefore their product C\_struct=W.c0\_cb*W.K\_call and e\_struct=W.epsilon\_MAIN have the asserted properties. \steptag{validated} \\[4pt]
\end{proofbox}

\end{document}
